\begin{table}[!h]
\centering\centering\centering
\caption{Effect on Test Scores}
\centering
\begin{threeparttable}
\fontsize{10}{12}\selectfont
\begin{tabular}[t]{lccccccc}
\toprule
\multicolumn{1}{c}{ } & \multicolumn{7}{c}{Experiment (Liberia)} \\
\cmidrule(l{3pt}r{3pt}){2-8}
  & (1) & (2) & (3) & (4) & (5) & (6) & (7)\\
\midrule
Peer Predicted Score & 0.163 &  &  &  & 0.095 &  & 0.096\\
 & (0.180) &  &  &  & (0.185) &  & (0.185)\\
Dispersion in Peer Scores &  & 0.992* &  &  &  & 0.373 & 0.038\\
 &  & (0.558) &  &  &  & (0.579) & (0.482)\\
Distance from Median Instruction &  &  & -0.357 &  & -0.045 & -0.227 & -0.045\\
 &  &  & (0.257) &  & (0.208) & (0.227) & (0.210)\\
Class Size &  &  &  & 0.004 & 0.001 & 0.002 & 0.001\\
 &  &  &  & (0.005) & (0.004) & (0.005) & (0.004)\\
\hline
\midrule
Number of Students & 3553 & 3553 & 3103 & 3905 & 3102 & 3102 & 3102\\
Number of Schools & 54 & 54 & 54 & 54 & 54 & 54 & 54\\
Number of Test Scores & 1 & 1 & 1 & 1 & 1 & 1 & 1\\
\bottomrule
\end{tabular}
\begin{tablenotes}[para]
\item Notes:
                      All specifications control for individual baseline test score,
                        grade dummies and randomization strata fixed effects. Each specification
                        also controls for the expected value of the respective independent variable
                         used in the specification. Midline test scores are used in place of endline test scores when the latter are unavailable. 
                         When both endline and midline test scores are available, the average of the two is used. 
                         Standard errors are clustered at the academy-grade
                          group level. ***, **, and * indicate significance at 1\%, 5\%, and 10\%.
\end{tablenotes}
\end{threeparttable}
\end{table}
